% !TEX root = ../main.tex
Lo studio si basa su modello del VTOL ottenuto attraverso la formulazione di Newton-Eulero a 6 DOF.

\section{Sistemi di Riferimento}
Per descrivere la dinamica del tricottero VTOL, si adottano due sistemi di riferimento principali, entrambi destrorsi:
\begin{itemize}
    \item \textbf{Sistema Inerziale (Global Frame) $\mathcal{F}_g$}: Fisso rispetto al terreno, definito dalla terna $(O_g, X_g, Y_g, Z_g)$. La posizione del velivolo è indicata dal vettore $\bm{P}_g = [x_g, y_g, z_g]^T$ e l'orientamento dagli angoli di Eulero $\bm{\Theta}_g = [\phi, \theta, \psi]^T$ (Roll, Pitch, Yaw).
    \item \textbf{Sistema Solidale (Body Frame) $\mathcal{F}_b$}: Con origine nel centro di massa (CoM) del velivolo e assi $(X_b, Y_b, Z_b)$ solidali alla struttura. Le velocità lineari e angolari sono definite rispettivamente come $\bm{V}_b = [u, v, w]^T$ e $\bm{\Omega}_b = [p, q, r]^T$.
\end{itemize}

\begin{figure}[h]
    \centering
    \includegraphics[width=0.95\textwidth]{sist_rif.png}
    \caption{Sistemi di riferimento.}
    \label{fig:sist_rif}
\end{figure}


\subsection{Matrici di Trasformazione}
La trasformazione delle grandezze vettoriali dal sistema solidale al corpo a quello inerziale avviene tramite la matrice di rotazione ortogonale $\mathbf{R}_{b \to g}$ (secondo la convenzione standard Z-Y-X).\\
\[
			R_x(\phi) =
			\begin{bmatrix}
				1 & 0 & 0 \\
				0 & \cos\phi & -\sin\phi \\
				0 & \sin\phi & \cos\phi
			\end{bmatrix}, \quad
			R_y(\theta) =
			\begin{bmatrix}
				\cos\theta & 0 & \sin\theta \\
				0 & 1 & 0 \\
				-\sin\theta & 0 & \cos\theta
			\end{bmatrix}\]
		
		\begin{equation}\quad
			R_z(\psi) =
			\begin{bmatrix}
				\cos\psi & -\sin\psi & 0 \\
				\sin\psi & \cos\psi & 0 \\
				0 & 0 & 1
			\end{bmatrix}
		\end{equation}

In forma esplicita, utilizzando la notazione $c_{\cdot}=\cos(\cdot)$ e $s_{\cdot}=\sin(\cdot)$:
\begin{equation}
    \mathbf{R}_{b \to g} = 
    \begin{bmatrix}
        c_\psi c_\theta & c_\psi s_\theta s_\phi - s_\psi c_\phi & c_\psi s_\theta c_\phi + s_\psi s_\phi \\
        s_\psi c_\theta & s_\psi s_\theta s_\phi + c_\psi c_\phi & s_\psi s_\theta c_\phi - c_\psi s_\phi \\
        -s_\theta       & c_\theta s_\phi                        & c_\theta c_\phi
    \end{bmatrix}
    \label{eq:matrix_R}
\end{equation}
Poiché ${R}_{b \to g}$ è una matrice ortogonale, è possibile passare dal sistema inerziale a quello solidale al corpo utilizzando la sua trasposta.

Si definisce adesso  $\mathbf{J}_{b \to g}$, matrice di trasformazione delle velocità angolari dal body frame al global frame:
\begin{equation}
    \mathbf{J}_{b \to g} = 
    \begin{bmatrix}
        1 & \sin\phi \tan\theta & \cos\phi \tan\theta \\
        0 & \cos\phi            & -\sin\phi \\
        0 & \frac{\sin\phi}{\cos\theta} & \frac{\cos\phi}{\cos\theta}
    \end{bmatrix}
    \label{eq:matrix_J}
\end{equation}
A differenza di $\mathbf{R}_{b \to g}$, la matrice $\mathbf{J}_{b \to g}$ presenta una singolarità matematica per $\theta = \pm 90^\circ$, nota come \textit{Gimbal Lock}. Questa relazione è cruciale per calcolare correttamente il tasso di virata ($\dot{\psi}$) partendo dal rateo di imbardata ($r$) misurato dai sensori.

Dunque, le equazioni definite nel body frame possono essere trasformate nel sistema di rifeirmento inerziale attraverso le sequenti relazioni:
\begin{equation}
    \bm{\dot{P}}_g = \mathbf{R}_{b \to g}(\bm{\Theta}_g) \bm{V}_b
    \label{eq:pos_global}
\end{equation}
\begin{equation}
    \bm{\dot{\Theta}}_g = \mathbf{J}_{b \to g}(\bm{\Theta}_g) \bm{\Omega}_b
    \label{eq:vel_ang_global}
\end{equation}

È possibile scrivere tali relazioni in forma matriciale:
\begin{equation}
        \begin{bmatrix} \bm{\dot{P}}_g \\ \bm{\dot{\Theta}}_g \end{bmatrix} = 
        \begin{bmatrix} \mathbf{R}_{b \to g}(\bm{\Theta}) & \mathbf{0} \\ \mathbf{0} & \mathbf{J}_{b \to g}(\bm{\Theta}) \end{bmatrix}
        \begin{bmatrix} \bm{V}_b \\ \bm{\Omega}_b \end{bmatrix}
\end{equation}

\clearpage

\section{Equazioni della Dinamica}
Il modello dinamico è derivato utilizzando la formulazione di Newton-Eulero per un corpo rigido a 6 gradi di libertà. Le equazioni vettoriali che descrivono il moto traslazionale e rotazionale sono:

\begin{equation}
    \begin{cases}
        m \bm{\dot{V}}_b + \bm{\Omega}_b \times (m \bm{V}_b) = \bm{F}_{\text{TOT}} \\
        \mathbf{I}_b \bm{\dot{\Omega}}_b + \bm{\Omega}_b \times (\mathbf{I}_b \bm{\Omega}_b) = \bm{M}_{\text{TOT}}
    \end{cases}
    \label{eq:newton_euler}
\end{equation}

Dove:
\begin{itemize}
    \item $m$: massa totale del tricottero;
    \item $\mathbf{I}$: matrice identità $3 \times 3$;
    \item $\mathbf{I}_b$: matrice d'inerzia del tricottero;
    \item $\mathbf{V}_b = [v_x, v_y, v_z]^\top$: velocità lineari espresse nel \textit{body frame};
    \item $\dot{\mathbf{V}}_b$: accelerazioni lineari espresse nel \textit{body frame};
    \item $\boldsymbol{\Omega}_b = [p, q, r]^\top$: velocità angolari (\textit{roll, pitch, yaw rate}) nel \textit{body frame};
    \item $\dot{\boldsymbol{\Omega}}_b$: accelerazioni angolari nel \textit{body frame};
    \item $\boldsymbol{\Omega}_b \times (m\mathbf{V}_b)$: termine di Coriolis;
    \item $\boldsymbol{\Omega}_b \times (\mathbf{I}_b\boldsymbol{\Omega}_b)$: termine giroscopico dovuto all'inerzia del tricottero;
    \item $\mathbf{F}_t$: forze totali agenti sul tricottero;
    \item $\mathbf{M}_t$: momenti totali applicati al tricottero.
\end{itemize}
\clearpage

\section{Analisi delle Forze}
Le forze agenti sul tricottero sono la risultante della spinta dei motori ($\bm{F}_{\text{TH}}$), della forza di gravità ($\bm{F}_g$) e delle forze aerodinamiche ($\bm{F}_{\text{AERO}}$):
\begin{equation}
    \bm{F}_{\text{TOT}} = \bm{F}_{\text{TH}} + \bm{F}_{GRAV} + \bm{F}_{\text{AERO}}
\end{equation}

\subsection{Configurazione Geometrica e Thrust}
Il velivolo presenta una configurazione a tre rotori asimmetrica; sono presenti due rotori posizionati anteriormente e uno posizionato posteriormente, secondo lo schema della figura \ref{fig:config_rot}:
\begin{figure}
    \centering
    \includegraphics[width=0.8\textwidth]{tricopter_2D.png}
    \caption{Configurazione dei rotori.}
    \label{fig:config_rot}
\end{figure}

I bracci dei rotori sono definiti nel sistema body dai vettori posizione $\bm{r}_i$:
\begin{itemize}
    \item $\bm{r}_{1,2} = [d_{mx}, \pm d_{my}, 0]^T$
    \item $\bm{r}_3 = [-d_{tx}, 0, 0]^T$
\end{itemize}

I rotori anteriori hanno un grado di libertà e possono inclinarsi attorno all'asse $Y_b$ rispettivamente di un angolo $\theta_1$ e $\theta_2$.
Il rotore di coda possiede due gradi di libertà, potendo ruotare attorno a $Y_b$ di un angolo $\theta_3$ e attorno a $Z_b$ di un angolo $\theta_4$.

Ogni rotore genera una spinta $f_i = k \omega_i^2$, dove $\omega_i$ è la velocità angolare del rotore $i$ e $k$ il coefficiente di spinta (\textit{thrust coefficient})[$N \cdot s^2/rad^2$]. Il vettore forza totale si ottiene sommando i contributi dei tre rotori ruotati dalle rispettive matrici di tilt.

% CORREZIONE: Uso di minipage al posto di subfigure per evitare errori di pacchetti
\begin{figure}[htbp]
    \centering
    % Prima immagine (sinistra) - [b] allinea alla base
    \begin{minipage}[b]{0.48\textwidth}
        \centering
        \vfill % Spinge l'immagine verso il basso se è più bassa dell'altra
        \includegraphics[width=\linewidth]{rotori_1_2.png}
        \vspace{0pt} % Punto di ancoraggio per l'allineamento
        \caption{Tilt rotori anteriori ($X_b, Z_b$)}
        \label{fig:rotori_ant}
    \end{minipage}
    \hfill 
    % Seconda immagine (destra) - [b] allinea alla base
    \begin{minipage}[b]{0.48\textwidth}
        \centering
        \vfill
        \includegraphics[width=\linewidth]{rotore_3.png}
        \vspace{0pt}
        \caption{Tilt rotore coda ($X_b, Y_b, Z_b$)}
        \label{fig:rotore_post}
    \end{minipage}
    % Nota: \label{fig:rotori_tilt} qui non funzionerà correttamente 
    % se non c'è una \caption generale della figura.
\end{figure}


Espandendo i termini di rotazione, la forza di spinta risultante $\bm{F}_{\text{TH}}$ assume la forma:
\begin{equation}
    \bm{F}_{\text{TH}} = 
    \begin{bmatrix}
    k\omega_1^2 c_{\theta_1} + k\omega_2^2 c_{\theta_2} + k\omega_3^2 c_{\theta_3} c_{\theta_4} \\
    k\omega_3^2 c_{\theta_3} s_{\theta_4} \\
    -k\omega_1^2 s_{\theta_1} - k\omega_2^2 s_{\theta_2} - k\omega_3^2 s_{\theta_3}
    \end{bmatrix}
    \label{eq:F_th}
\end{equation}

Inizialmente il design del drone prevedeva che i rotori anteriori fossero in linea con le ali; tuttavia questa configurazione è stata superata 
per rispondere a dei requisiti meccanici fondamentiali:
\begin{itemize}
    \item \textbf{Ottimizzazione del Baricentro (CG):} I motori e i meccanismi di \textit{tilting} rappresentano le masse concentrate maggiori. Avanzarli permette di posizionare il $CG$ del velivolo davanti al Punto Neutro dell'ala. Questa condizione è indispensabile per la stabilità statica longitudinale in volo di crociera, garantendo che il drone sia naturalmente autostabilizzante rispetto alle perturbazioni di beccheggio.
    \item \textbf{Braccio di Leva:} Spostando i rotori anteriori longitudinalmente rispetto al $CG$, si incrementa il braccio di leva per il controllo del beccheggio. Per un tricottero, ciò permette una gestione più efficiente del momento longitudinale, riducendo il carico di lavoro del rotore di coda e prevenendo la saturazione degli attuatori durante le manovre correttive.
\end{itemize}


\subsection{Gravità e Aerodinamica}
La forza peso viene ruotata nel sistema solidale al corpo tramite la trasposta della matrice di rotazione:
\begin{equation}
    \bm{F}_{GRAV} = \mathbf{R}_{b \to g}^T \begin{bmatrix} 0 \\ 0 \\ mg \end{bmatrix} = mg \begin{bmatrix} -\sin\theta \\ \cos\theta \sin\phi \\ \cos\theta \cos\phi \end{bmatrix}
    \label{eq:F_grav}
\end{equation}

Per una modellazione aerodinamica più fedele al comportamento in fase di crociera, si definiscono l'angolo di attacco $\alpha$ e l'angolo di scivolamento laterale $\beta$ in funzione delle velocità nel body frame:
\begin{equation}
    \alpha = \arctan\left(\frac{w}{u}\right), \quad \beta = \arcsin\left(\frac{v}{V_{tas}}\right)
\end{equation}

Le forze aerodinamiche includono il contributo delle ali ($\bm{F}_{\text{{wing}}}$) e la resistenza della fusoliera ($\bm{D}_{\text{body}}$):
\begin{equation}
    \bm{F}_{\text{AERO}} = \bm{F}_{\text{{wing}}} + \bm{D}_{\text{body}}
\end{equation}

La resistenza del corpo è modellata lungo gli assi principali in funzione del segno della velocità ($\text{sgn}(v)$) per garantire coerenza fisica:
\begin{equation}
    \bm{D}_{\text{body}} = -\frac{1}{2} \rho 
    \begin{bmatrix}
        \text{sgn}(u) u^2 S_{x} C_{Dx} \\
        \text{sgn}(v) v^2 S_{y} C_{Dy} \\
        \text{sgn}(w) w^2 S_{z} C_{Dz}
    \end{bmatrix}
\end{equation}

Per le superfici alari, si considerano Portanza ($L$) e Resistenza ($D$):
\begin{equation}
	 \mathbf{F}_{\text{wing}} =
	 \begin{bmatrix}
	 	F_{D_1} + F_{D_2} \\
	 	0 \\
	 	F_{L_1} + F_{L_2} + F_{D_{z1}} + F_{D_{z2}}
	 \end{bmatrix}
\end{equation}

dove si è considerata nulla (per simmetria) la componente della forza esercitata lungo l'asse $Y_{b}$ e con: 
\begin{itemize}
    \item $F_{D_i}$ : forza di resistenza (drag) lungo l'asse $X_{b}$ generata dalla superficie alare $i$-esima
    \item $F_{L_i}$ : forza di portanza (lift) generata dalla superficie alare $i$-esima
    \item $F_{D_{zi}}$ : forza di resistenza (drag) lungo l'asse $Z_{b}$ generata dalla superficie alare $i$	-esima
\end{itemize}

\begin{figure}[h]
    \centering
    \includegraphics[width=0.7\textwidth]{forze_x.jpg}
    \caption{Portanza e Resistenza lungo l'asse $X_b$}
\end{figure}

\begin{equation}
    \bm{F}_{\text{{wing}}} = 
    \begin{bmatrix}
        -\frac{1}{2}\rho S C_D(u) u |u| \\
        0 \\
        -\frac{1}{2}\rho S C_L(u) u^2 - -\frac{1}{2}\rho S C_{Dz} w^2
    \end{bmatrix}
\end{equation}

dove:
\begin{itemize}
    \item $u$ : velocità [$m/s$] lungo l'asse $X_{B}$
    \item $w$ : velocità [$m/s$] lungo l'asse $Z_{B}$
    \item $\rho$ : densità dell'aria [$kg/m^3$]
    \item $S$ : superficie alare  [$m^2$]
    \item $C_D$ : coefficiente di resistenza aerodinamica lungo l'asse $X_{b}$
    \item $C_L$ : coefficiente di portanza 
    \item $C_{D_{z}}$ : coefficiente di resistenza aerodinamica lungo l'asse $Z_{b}$
\end{itemize}

\begin{figure}[H]
    \centering
    \includegraphics[width=0.7\textwidth]{forze_z.jpg}
    \caption{Resistenza lungo l'asse $Z_b$}
\end{figure}

\begin{equation*}
	  	\mathbf{F}_{TOT} = \mathbf{F}_{TH} + \mathbf{F}_{GRAV} + \mathbf{F}_{AERO} = 
	  	\end{equation*}
	  	
	  	 \begin{equation*}
	  		=  
	  		\begin{bmatrix}
	  			\cos(\theta_{1}) & \cos(\theta_{2}) & {\cos(\theta_{3})\cos(\theta_{4})} \\
	  			0                 & 0                 & {\cos(\theta_{3})\sin(\theta_{4})} \\
	  			-\sin(\theta_{1}) & -\sin(\theta_{2}) & {-\sin(\theta_{3})}
	  		\end{bmatrix}
	  		\begin{bmatrix}
	  			k\omega_{1}^2 \\
	  			k\omega_{2}^2 \\
	  			k\omega_{3}^2 
	  		\end{bmatrix} +  \begin{bmatrix}
	  		-\sin(\theta) \\
	  		\cos(\theta)\sin(\phi) \\
	  		\cos(\theta)\cos(\phi) 
	  		\end{bmatrix} m g +  
	  	\end{equation*}
	  \begin{equation}
	  	+
	   \begin{bmatrix}
	  	-\operatorname{sgn}(v_x) \,\frac{1}{2}\rho\, v_x^2\, s\, (C_d + C_d) \\
	  	0 \\
	  	-\frac{1}{2}\rho\, v_x^2\, s\, (C_l + C_l) - \operatorname{sgn}(v_z)\,\frac{1}{2}\rho\, v_z^2\, s\, (C_{d,z} + C_{d,z})
	  \end{bmatrix}
	\label{eq:f_tot}
\end{equation}

\clearpage

\section{Analisi dei Momenti}
Dall'equazione della dinamica si ottiene che :
	  
\begin{equation}
I_B \, \dot{\boldsymbol{\Omega}}_B +
\underbrace{\boldsymbol{\Omega}_B \times (I_B \boldsymbol{\Omega}_B)}_{\text{Momento giroscopico del corpo (tricottero)}}
= \mathbf{M}_{\text{TOT}}
\label{eq:eq_momenti}
\end{equation}

\vspace{0.3cm}
dove la matrice di inerzia del tricottero è pari a:

\begin{equation}
    \mathbf{I}_{\text{B}} =
    \begin{bmatrix}
        I_{\text{xx}} & 0 & 0 \\
        0 & I_{\text{yy}} & 0 \\
        0 & 0 & I_{\text{zz}}
    \end{bmatrix}
\end{equation}

\vspace{0.3cm}
Di conseguenza si ottiene che:

\begin{equation}
    I_B \, \dot{\boldsymbol{\Omega}}_B =
    \mathbf{M}_{\text{TOT}} -
    \boldsymbol{\Omega}_B \times (I_B \boldsymbol{\Omega}_B)
    \label{eq:eq_momenti_esplicita}
\end{equation}


Il momento totale agente sul baricentro è la somma vettoriale delle seguenti componenti:
\begin{equation}
    \bm{M}_{\text{TOT}} = \bm{M}_{\text{TH}} + \bm{M}_{\text{TOR}} + \bm{M}_{\text{AERO}} + \bm{M}_{\text{GIRO}} + \bm{M}_{\text{PINNA}}
    \label{eq:M_tot}
\end{equation}

\subsection{Momento di Thrust e Torcente}
Il contributo principale $\bm{M}_{\text{TH}}$ deriva dal prodotto vettoriale tra i bracci $\bm{r}_i$ e le forze di spinta. Utilizzando le coordinate geometriche definite in precedenza:
\begin{equation}
    \bm{M}_{\text{TH}} = \sum_{i=1}^{3} (\bm{r}_i \times \bm{F}_{\text{TH},i}) =
    \begin{bmatrix} 
    -d_{my} (k\omega_1^2 s_{\theta_1}) + d_{my} (k\omega_2^2 s_{\theta_2}) \\ 
    d_{mx} (k\omega_1^2 s_{\theta_1}) + d_{mx} (k\omega_2^2 s_{\theta_2}) + d_{tx} (k\omega_3^2 s_{\theta_3}) \\ 
    -d_{my} (k\omega_1^2 c_{\theta_1}) + d_{my} (k\omega_2^2 c_{\theta_2}) + d_{tx} (k\omega_3^2 c_{\theta_3} s_{\theta_4})
    \end{bmatrix}
    \label{eq:M_th}
\end{equation}

Inoltre, la rotazione delle eliche genera un momento torcente resistivo ($\tau_i = b \omega_i^2$). Poiché i rotori sono inclinabili, questo momento deve essere proiettato sugli assi del {body} frame:
\begin{equation}
    \bm{M}_{\text{TOR}} = 
    \begin{bmatrix} 
    b\omega_1^2 c_{\theta_1} - b\omega_2^2 c_{\theta_2} + b\omega_3^2 c_{\theta_3} c_{\theta_4} \\ 
    b\omega_3^2 c_{\theta_3} s_{\theta_4} \\ 
    -b\omega_1^2 s_{\theta_1} + b\omega_2^2 s_{\theta_2} - b\omega_3^2 s_{\theta_3}
    \end{bmatrix}
    \label{eq:M_tor}
\end{equation}

\subsection{Effetti Aerodinamici e Giroscopici}
\textbf{Momento Aerodinamico ($\bm{M}_{\text{AERO}}$):} Include le coppie generate dal disallineamento tra il Centro di Pressione (CP) delle ali e il CoM, oltre ai momenti di resistenza aerodinamica della fusoliera.
\begin{equation}
    \bm{M}_{AERO} = 
    \begin{bmatrix}
        (l_{sx_{y}}+l_{dx_{y}})\left(-\tfrac{C_l \rho s v_x^2}{2} - \tfrac{C_{d_z} \rho s v_z^2}{2}\,\text{sign}(v_z)\right) \\[10pt]
        -\Big[ (l_{sx_{x}}+l_{dx_{x}})\left(-\tfrac{C_l \rho s v_x^2}{2} - \tfrac{C_{d_z} \rho s v_z^2}{2}\,\text{sign}(v_z)\right) 
        - (l_{sx_{z}}+l_{dx_{z}})\left(-\tfrac{C_d \rho s v_x^2}{2}\,\text{sign}(v_x)\right)\Big] \\[10pt]
        (l_{sx_{y}}+l_{dx_{y}})\,\tfrac{C_d \rho s v_x^2}{2}\,\text{sign}(v_x)
    \end{bmatrix}
    \label{eq:M_aero}
\end{equation}

\textbf{Momento Giroscopico ($\bm{M}_{\text{GIRO}}$):} Dovuto all'interazione tra la velocità angolare di rotazione dell'elica $\omega_i$ e la velocità di tilting dei servomotori ($\bm{\omega}_{\text{tilt}}$). Questo termine è critico durante le transizioni rapide:
\begin{equation}
    \bm{M}_{\text{GIRO}}= \sum_{i=1}^{3} \mathbf{J}_{rot} (\bm{\omega}_{\text{tilt},i} \times \bm{\omega}_{i})
\end{equation}
\begin{equation}
    \mathbf{M}_{\text{GIRO}} = 
    \resizebox{0.93\textwidth}{!}{$
        \begin{bmatrix}
            -I_{\text{rot}_{12_{z}}}\,\omega_{\text{tilt}_{1}} \sin(\theta_{1}) \omega_{1}
            + I_{\text{rot}_{12_{z}}}\,\omega_{\text{tilt}_{2}} \sin(\theta_{2})\omega_{2}
            - I_{\text{rot}_{3_{z}}}\,\omega_{\text{tilt}_{3,xz}} \sin(\theta_{3}) \omega_{3}
            - I_{\text{rot}_{3_{y}}}\,\omega_{\text{tilt}_{3,xy}} \cos(\theta_{3}) \sin(\theta_{4}) \omega_{3} \\[6pt]
            + I_{\text{rot}_{3_{x}}}\,\omega_{\text{tilt}_{3,xy}} \cos(\theta_{3}) \cos(\theta_{4})\omega_{3} \\[6pt]
            -I_{\text{rot}_{12_{x}}}\,\omega_{\text{tilt}_{1}} \cos(\theta_{1}) \omega_{1}
            + I_{\text{rot}_{12_{x}}}\,\omega_{\text{tilt}_{2}} \cos(\theta_{2}) \omega_{2}
            - I_{\text{rot}_{3_{x}}}\,\omega_{\text{tilt}_{3,xz}} \cos(\theta_{3}) \cos(\theta_{4}) \omega_{3}
        \end{bmatrix}
        $}
        \label{eq:M_giro}
\end{equation}

Il modello è caratterizzato da un secondo termine giroscopico, dovuto alla rotazione del tricottero nel body frame:
\begin{equation*}
        \boldsymbol{M}_{body} =
        \boldsymbol{\Omega}_{B} \times (I_B \boldsymbol{\Omega}_{B}) =
            \begin{bmatrix}
            p \\
            q \\
            r
        \end{bmatrix}
        \times
        \begin{bmatrix}
            I_{\text{xx}} & 0 & 0 \\
            0 & I_{\text{yy}} & 0 \\
            0 & 0 & I_{\text{zz}}
        \end{bmatrix}
        \begin{bmatrix}
            p\\
            q\\
            r
        \end{bmatrix}
        =
        \end{equation*}
        
        \begin{equation}
        =
        \begin{bmatrix}
            (I_{zz}-I_{yy})	qr \\
            (I_{xx}-I_{zz})	pr\\
            (I_{yy}-I_{xx})	pq
        \end{bmatrix}
        \label{eq:M_body}
        \end{equation}

\textbf{Momento della Pinna ($\bm{M}_{\text{PINNA}}$):} La superficie verticale di coda garantisce la stabilità direzionale (\textit{Weathercock stability}) e lo smorzamento dell'imbardata.
\begin{equation}
    \mathbf{M}_{\text{PINNA}} = 
    \begin{bmatrix}
    M_{\text{pinna},x} \\ M_{\text{pinna},y} \\ M_{\text{pinna},z}
    \end{bmatrix} =
    \begin{bmatrix}
    -p\alpha_{0x}\ - \alpha_{1x}\,v_y^2 \\
    0 \\
    -r\left( \alpha_{0z}\ + \alpha_{1z}\,v_y^2 \right)
    \end{bmatrix}
    \label{eq:M_pinna}
\end{equation}

La presenza della pinna è stata dibattuta in fase di revisione del modello, 
tuttavia, si è stabilito il suo inserimento poiché garantisce una maggiore 
stabilità di imbardata soprattutto in considerazione della coppia
di reazione residua prodotta dal rotore di coda.  

\subsection{Momento totale} 
Utilizzando i risultati ottenuti sopra, si ottiene l'espressione del momento totale:

	    \begin{equation*}
	    	I_B \, \dot{\boldsymbol{\Omega}}_B =
	    	\mathbf{M}_{\text{TOT}} -
	    	\boldsymbol{\Omega}_B \times (I_B \boldsymbol{\Omega}_B) = \mathbf{M}_{\text{TH}} + \mathbf{M}_{\text{TOR}}+\mathbf{M}_{\text{AERO}}  +  \mathbf{M}_{\text{GIRO}} +  \mathbf{M}_{\text{PINNA}}  -
	    	\boldsymbol{\Omega}_B \times (I_B \boldsymbol{\Omega}_B) =
	    \end{equation*}
	    
	     \begin{equation*}
	    	 = \resizebox{0.99\textwidth}{!}{$
	    	 	\begin{bmatrix}
	    	 		d_{my} \, k \omega_2^2 \sin(\theta_2) 
	    	 		- d_{my} \, k \omega_1^2 \sin(\theta_1) 
	    	 		- d_{ty} \, k \omega_3^2 \sin(\theta_3) 
	    	 		- d_{tz} \, k \omega_3^2 \cos(\theta_3) \sin(\theta_4) \\[6pt]
	    	 		d_{mz} \, k \omega_1^2 \cos(\theta_1) 
	    	 		+ d_{mz} \, k \omega_2^2 \cos(\theta_2) 
	    	 		+ d_{mx} \, k \omega_1^2 \sin(\theta_1) 
	    	 		+ d_{mx} \, k \omega_2^2 \sin(\theta_2) 
	    	 		+ d_{tx} \, k \omega_3^2 \sin(\theta_3) 
	    	 		+ d_{tz} \, k \omega_3^2 \cos(\theta_3) \cos(\theta_4) \\[6pt]
	    	 		d_{my} \, k \omega_2^2 \cos(\theta_2) 
	    	 		- d_{my} \, k \omega_1^2 \cos(\theta_1) 
	    	 		- d_{ty} \, k \omega_3^2 \cos(\theta_3) \cos(\theta_4) 
	    	 		+ d_{tx} \, k \omega_3^2 \cos(\theta_3) \sin(\theta_4)
	    	 	\end{bmatrix}
	    	 	$} +
	    \end{equation*}
	    
	   	  \begin{equation*}
	  	+	\begin{bmatrix}
	  		(b\omega_{1}^2+I_{\text{rot}_{12_{z}}}\dot\omega_{1}) \cos(\theta_1) -	(b\omega_{2}^2+I_{\text{rot}_{12_{z}}}\dot\omega_{2}) \cos(\theta_2) + (b\omega_{3}^2+I_{\text{rot}_{3_{z}}}\dot\omega_{3})\cos(\theta_3)\cos(\theta_4) \\
	  		(b\omega_{3}^2+I_{\text{rot}_{3_{z}}}\dot\omega_{3})\cos(\theta_3)\sin(\theta_4) \\
	  		-	(b\omega_{1}^2+I_{\text{rot}_{12_{z}}}\dot\omega_{1})\sin(\theta_1)+ 	(b\omega_{2}^2+I_{\text{rot}_{12_{z}}}\dot\omega_{2})  \sin(\theta_2)	-	(b\omega_{3}^2+I_{\text{rot}_{3_{z}}}\dot\omega_{3})\sin(\theta_3)
	  	\end{bmatrix}+
	  \end{equation*}
	    
	    \begin{equation*}
	    	+
	    	\begin{bmatrix}
	    		(l_{sx_{y}}+l_{dx_{y}})\left(-\tfrac{C_l \rho s v_x^2}{2} - \tfrac{C_{d_z} \rho s v_z^2}{2}\,\text{sign}(v_z)\right) \\[10pt]
	    		-\Big[ (l_{sx_{x}}+l_{dx_{x}})\left(-\tfrac{C_l \rho s v_x^2}{2} - \tfrac{C_{d_z} \rho s v_z^2}{2}\,\text{sign}(v_z)\right) 
	    		- (l_{sx_{z}}+l_{dx_{z}})\left(-\tfrac{C_d \rho s v_x^2}{2}\,\text{sign}(v_x)\right)\Big] \\[10pt]
	    		(l_{sx_{y}}+l_{dx_{y}})\,\tfrac{C_d \rho s v_x^2}{2}\,\text{sign}(v_x)
	    	\end{bmatrix}
	    	+
	    \end{equation*}
	    
	     \begin{equation*}
	    	+
	    	 \resizebox{0.93\textwidth}{!}{$
	    		\begin{bmatrix}
	    			-I_{\text{rot}_{12_{z}}}\,\omega_{\text{tilt}_{1}} \sin(\theta_{1}) \omega_{1}
	    			+ I_{\text{rot}_{12_{z}}}\,\omega_{\text{tilt}_{2}} \sin(\theta_{2})\omega_{2}
	    			- I_{\text{rot}_{3_{z}}}\,\omega_{\text{tilt}_{3,xz}} \sin(\theta_{3}) \omega_{3}
	    			- I_{\text{rot}_{3_{y}}}\,\omega_{\text{tilt}_{3,xy}} \cos(\theta_{3}) \sin(\theta_{4}) \omega_{3} \\[6pt]
	    			+ I_{\text{rot}_{3_{x}}}\,\omega_{\text{tilt}_{3,xy}} \cos(\theta_{3}) \cos(\theta_{4})\omega_{3} \\[6pt]
	    			-I_{\text{rot}_{12_{x}}}\,\omega_{\text{tilt}_{1}} \cos(\theta_{1}) \omega_{1}
	    			+ I_{\text{rot}_{12_{x}}}\,\omega_{\text{tilt}_{2}} \cos(\theta_{2}) \omega_{2}
	    			- I_{\text{rot}_{3_{x}}}\,\omega_{\text{tilt}_{3,xz}} \cos(\theta_{3}) \cos(\theta_{4}) \omega_{3}
	    		\end{bmatrix}
	    		$}
	    	+
	    \end{equation*}
	    
	    \begin{equation}
	    	+
	    	\begin{bmatrix}
	    		-p\alpha_{0x}\ - \alpha_{1x}\,v_y^2 \\
	    		0 \\
	    		-r\left( \alpha_{0z}\ + \alpha_{1z}\,v_y^2 \right)
	    	\end{bmatrix}
	    	- \begin{bmatrix}
	    		(I_{zz}-I_{yy})	qr \\
	    		(I_{xx}-I_{zz})	pr\\
	    		(I_{yy}-I_{xx})	pq
	    	\end{bmatrix}
            \label{eq:matrix_M_tot}
	    \end{equation}
\clearpage
